\subsection{Feature Selection Metrics}

To narrow the graph feature set, we audit each candidate feature with a held-out AUC protocol instead of relying on intuition alone. The cached edge table is joined with node-level statistics, filtered to remove self-loops and user-edge metadata, and labeled by whether each source-destination pair appears in the LANL red-team ground truth. A stratified 50/50 split reserves one half for calibrating the feature ranking and the other half for checking generalization. For each feature, we compute ROC AUC using that feature alone; if red-team edges have a lower mean than benign edges, we invert the AUC so the score measures separability regardless of direction. This puts features like low destination degree and high NTLM usage on the same scale.

The audit also records supporting statistics for each feature: number of unique values, variance, red-team mean, benign mean, and the difference between those means. These extra checks matter because binary authentication indicators can have low variance yet still separate attacks well, so a PCA-style variance ranking would miss useful signals. In the current audit, the strongest single features are \texttt{is\_ntlm} (AUC 0.933), \texttt{source\_fan\_out} (AUC 0.906), \texttt{dst\_in\_degree} (AUC 0.819), \texttt{is\_network\_logon} (AUC 0.817), and \texttt{dst\_fan\_out\_ratio} (AUC 0.817). The pattern is consistent with lateral movement: remote logons, NTLM-heavy authentication, unusual outbound spread from a source, and destinations that do not behave like ordinary high-traffic services.

For the final scoring model, we use held-out AUC as a screening metric, not as an automatic feature selector. Features above the current AUC threshold of 0.6 are reasonable candidates for the graph score, while weak or constant features such as \texttt{temporal\_decay\_weight}, \texttt{is\_unusual\_dst\_port}, and \texttt{duration\_zscore} should stay out unless later runs show stronger evidence. This gives us a smaller and more interpretable set while still keeping signals from authentication semantics, graph structure, and timing or traffic behavior. The main caveat is that the audit still depends on LANL red-team labels, so the chosen feature set should be checked again on held-out windows or another dataset before we treat it as final.
